% GENERATED FILE. Edit canonical lessons, then run npm run generate:books.
% canonical-source-hash: fnv1a64:de2df28170f8e7c6
% canonical-lessons: PT-C461-lagarto, PT-C461-vibora, PT-C461-sapo, PT-C461-caracol, PT-C461-rastejar

\chapter{Bichos pequenos}
\label{ch:pt-bichos-pequenos}
\hlchaptermodality{461}

\begin{chapteropening}
\textbf{By the end of this chapter:} \emph{I can name small creatures on a walk: the lizard, the viper, the toad, the snail, and crawling.}

Complete the last lesson of chapter 461: I can name small creatures on a walk: the lizard, the viper, the toad, the snail, and crawling.
\end{chapteropening}

\section[o lagarto]{o lagarto — the lizard}

\label{lesson:PT-C461-lagarto}



\begin{quote}
Before the new one: say the Portuguese for the falcon, then the Portuguese for to graze.
\end{quote}

\subsection*{You'll want to know: o lagarto}
\textbf{o lagarto} — \textquotedblleft{}the lizard\textquotedblright{}.

\begin{quote}
Um lagarto apanha sol numa pedra.
\end{quote}

\begin{quote}
\emph{A lizard is sunbathing on a stone.}
\end{quote}

A small one is \emph{a lagartixa}.

\begin{cousinweb}
\emph{lagarto} is inherited from Vulgar Latin \emph{lacartus}, from Latin \emph{lacertus} 'lizard' (EtymDB).
\end{cousinweb}

\begin{grammarlens}[title={What you've built}]
One more word for telling, explaining and describing at B1.
\end{grammarlens}

\subsection*{Your turn}
\begin{itemize}
\raggedright
  \item \emph{Say it:} \emph{o lagarto}
  \item \emph{Say it:} \emph{o lagarto}, once more
  \item \emph{Say it:} a sentence of your own with \emph{o lagarto}
  \item \emph{Recall:} say the Portuguese for the falcon, then the Portuguese for to graze, then say \emph{o lagarto} again
\end{itemize}

\begin{tcolorbox}[breakable,colback=teal!4,colframe=teal!35!black,title={Before you move on}]
What does o lagarto mean? (\textquotedblleft{}The lizard\textquotedblright{}.) Say it once more.
\end{tcolorbox}
\section[a víbora]{a víbora — the viper, the adder}

\label{lesson:PT-C461-vibora}



\begin{quote}
Before the new one: say the Portuguese for to graze, then the Portuguese for the lizard.
\end{quote}

\subsection*{You'll want to know: a víbora}
\textbf{a víbora} — \textquotedblleft{}the viper, the adder\textquotedblright{}.

\begin{quote}
Na serra há víboras, por isso andamos com cuidado.
\end{quote}

\begin{quote}
\emph{There are vipers in the hills, so we walk carefully.}
\end{quote}

Figurative: a spiteful person.

\begin{cousinweb}
\emph{víbora} comes from Latin \emph{vīpera} 'viper' (EtymDB).
\end{cousinweb}

\begin{grammarlens}[title={What you've built}]
One more word for telling, explaining and describing at B1.
\end{grammarlens}

\subsection*{Your turn}
\begin{itemize}
\raggedright
  \item \emph{Say it:} \emph{a víbora}
  \item \emph{Say it:} \emph{a víbora}, once more
  \item \emph{Say it:} a sentence of your own with \emph{a víbora}
  \item \emph{Recall:} say the Portuguese for to graze, then the Portuguese for the lizard, then say \emph{a víbora} again
\end{itemize}

\begin{tcolorbox}[breakable,colback=teal!4,colframe=teal!35!black,title={Before you move on}]
What does a víbora mean? (\textquotedblleft{}The viper, the adder\textquotedblright{}.) Say it once more.
\end{tcolorbox}
\section[o sapo]{o sapo — the toad}

\label{lesson:PT-C461-sapo}



\begin{quote}
Before the new one: say the Portuguese for the lizard, then the Portuguese for the viper.
\end{quote}

\subsection*{You'll want to know: o sapo}
\textbf{o sapo} — \textquotedblleft{}the toad\textquotedblright{}.

\begin{quote}
À noite ouvimos os sapos no lago.
\end{quote}

\begin{quote}
\emph{At night we hear the toads in the lake.}
\end{quote}

\emph{Engolir um sapo} — to swallow a toad: to put up with something unpleasant.

\begin{cousinweb}
The origin of \emph{sapo} is not settled; it may be a pre-Roman word of the Iberian Peninsula, and EtymDB lists it without a source.
\end{cousinweb}

\begin{grammarlens}[title={What you've built}]
One more word for telling, explaining and describing at B1.
\end{grammarlens}

\subsection*{Your turn}
\begin{itemize}
\raggedright
  \item \emph{Say it:} \emph{o sapo}
  \item \emph{Say it:} \emph{o sapo}, once more
  \item \emph{Say it:} a sentence of your own with \emph{o sapo}
  \item \emph{Recall:} say the Portuguese for the lizard, then the Portuguese for the viper, then say \emph{o sapo} again
\end{itemize}

\begin{tcolorbox}[breakable,colback=teal!4,colframe=teal!35!black,title={Before you move on}]
What does o sapo mean? (\textquotedblleft{}The toad\textquotedblright{}.) Say it once more.
\end{tcolorbox}
\section[o caracol]{o caracol — the snail; the curl}

\label{lesson:PT-C461-caracol}



\begin{quote}
Before the new one: say the Portuguese for the viper, then the Portuguese for the toad.
\end{quote}

\subsection*{You'll want to know: o caracol}
\textbf{o caracol} — \textquotedblleft{}the snail; the curl\textquotedblright{}.

\begin{quote}
Depois da chuva, há caracóis no jardim.
\end{quote}

\begin{quote}
\emph{After the rain, there are snails in the garden.}
\end{quote}

In Portugal, \emph{caracóis} are a summer snack. \emph{Escada de caracol} — a spiral staircase.

\begin{cousinweb}
The origin of \emph{caracol} is not settled, and EtymDB has no entry for it.
\end{cousinweb}

\begin{grammarlens}[title={What you've built}]
One more word for telling, explaining and describing at B1.
\end{grammarlens}

\subsection*{Your turn}
\begin{itemize}
\raggedright
  \item \emph{Say it:} \emph{o caracol}
  \item \emph{Say it:} \emph{o caracol}, once more
  \item \emph{Say it:} a sentence of your own with \emph{o caracol}
  \item \emph{Recall:} say the Portuguese for the viper, then the Portuguese for the toad, then say \emph{o caracol} again
\end{itemize}

\begin{tcolorbox}[breakable,colback=teal!4,colframe=teal!35!black,title={Before you move on}]
What does o caracol mean? (\textquotedblleft{}The snail; the curl\textquotedblright{}.) Say it once more.
\end{tcolorbox}
\section[rastejar]{rastejar — to crawl, to creep}

\label{lesson:PT-C461-rastejar}



\begin{quote}
Before the new one: say the Portuguese for the toad, then the Portuguese for the snail.
\end{quote}

\subsection*{You'll want to know: rastejar}
\textbf{rastejar} — \textquotedblleft{}to crawl, to creep\textquotedblright{}.

\begin{quote}
As cobras rastejam pelo chão.
\end{quote}

\begin{quote}
\emph{Snakes crawl along the ground.}
\end{quote}

Built on \emph{o rasto}, the trail.

\begin{cousinweb}
\emph{rastejar} is built on \emph{rasto} 'trail, track', which is usually traced to Latin \emph{rāstrum} 'rake' (Lewis \& Short), itself linked with \emph{rādere} 'to scrape': the marks a rake leaves on the ground.
\end{cousinweb}

\begin{grammarlens}[title={What you've built}]
One more word for telling, explaining and describing at B1.
\end{grammarlens}

\subsection*{Your turn}
\begin{itemize}
\raggedright
  \item \emph{Say it:} \emph{rastejar}
  \item \emph{Say it:} \emph{rastejar}, once more
  \item \emph{Say it:} a sentence of your own with \emph{rastejar}
  \item \emph{Recall:} say the Portuguese for the toad, then the Portuguese for the snail, then say \emph{rastejar} again
\end{itemize}

\begin{tcolorbox}[breakable,colback=teal!4,colframe=teal!35!black,title={Before you move on}]
What does rastejar mean? (\textquotedblleft{}To crawl, to creep\textquotedblright{}.) Say it once more.
\end{tcolorbox}
